\documentclass[../general/general]{subfiles}

\begin{document}

\chapter{Лекция №12}
\noindent\textit{А.\,М.~Белемук \hfill 8 мая 2021}
\vspace{1em}
\ifSubfilesClassLoaded{\tableofcontents\newpage}{}

Сегодня мы продолжаем теорию слабо неидеального бозе-газа с того места, на котором остановились. Напомним: мы хотели вычислить коммутатор $[\hat{a}_\mathbf{k},\,\hat{V}]$ и написать уравнение движения для $\hat{a}_\mathbf{k}$ в представлении Гейзенберга. Первый коммутатор $[\hat{a}_\mathbf{k},\hat{H}_0] = \xi_\mathbf{k}\hat{a}_\mathbf{k}$ мы уже получили. Теперь беремся за второй.


% =====================================================================
\section{Завершение вычисления $[\hat{a}_\mathbf{k},\,\hat{V}]$}
% =====================================================================

Оператор взаимодействия в импульсном представлении:
\begin{equation}
    \hat{V} = \frac{v_0}{2V}\sum_{\mathbf{p},\mathbf{k}',\mathbf{q}}
    \hat{a}^\dagger_{\mathbf{p}+\mathbf{q}}\hat{a}^\dagger_{\mathbf{k}'-\mathbf{q}}\hat{a}_{\mathbf{k}'}\hat{a}_\mathbf{p}
    \label{eq:V}
\end{equation}

Нам нужно $[\hat{a}_\mathbf{k},\,\hat{a}^\dagger_1\hat{a}^\dagger_2\hat{a}_3\hat{a}_4]$. Раскрываем коммутатор через линейность:
\begin{equation}
    [\hat{a},BCDE] = [\hat{a},B]CDE + B[\hat{a},C]DE + BC[\hat{a},D]E + BCD[\hat{a},E]
\end{equation}
Оператор $\hat{a}_\mathbf{k}$ коммутирует только с $\hat{a}^\dagger_\mathbf{k}$ (даёт $\delta_{\mathbf{k}\mathbf{k}'}$) и с операторами уничтожения ($[\hat{a},\hat{a}]=0$). Поэтому ненулевой вклад дают только слагаемые, содержащие $[\hat{a}_\mathbf{k},\hat{a}^\dagger_{\cdots}]$. Из четырёх операторов в $\hat{V}$ первые два — операторы рождения, вторые два — уничтожения. Вклад дают коммутации с первым и вторым операторами рождения:

\begin{align}
    [\hat{a}_\mathbf{k},\,\hat{V}] &= \frac{v_0}{2V}\sum_{\mathbf{p},\mathbf{q}}
    \Bigl(\delta_{\mathbf{k},\mathbf{p}+\mathbf{q}}\,\hat{a}^\dagger_{\mathbf{k}-\mathbf{q}}\hat{a}_{\mathbf{k}-\mathbf{q}}\hat{a}_\mathbf{p}
    + \delta_{\mathbf{k},\mathbf{k}'-\mathbf{q}}\,\hat{a}^\dagger_{\mathbf{p}+\mathbf{q}}\hat{a}_{\mathbf{k}'}\hat{a}_\mathbf{p}\Bigr)
\end{align}

После замены переменных суммирования оба слагаемых дают одинаковый вклад, и результат:
\begin{equation}
    [\hat{a}_\mathbf{k},\,\hat{V}] = \frac{v_0}{V}\sum_{\mathbf{p},\mathbf{q}}
    \hat{a}^\dagger_{\mathbf{p}+\mathbf{q}}\hat{a}_{\mathbf{k}+\mathbf{q}}\hat{a}_\mathbf{p}
    \label{eq:comm_V_exact}
\end{equation}

Это точное выражение, приближений пока нет.

\subsection{Приближение конденсата}

Теперь делаем главное приближение. При $T \approx 0$ в конденсате сосредоточена бо́льшая часть частиц ($N_0 \sim N$). Доминирующие процессы в~\eqref{eq:comm_V_exact} — те, в которых два из трёх операторов уничтожают или рождают частицы с импульсом $\mathbf{p} = 0$ (конденсатные частицы). То есть мы оставляем только \emph{линейные} по наднконденсатным операторам слагаемые: либо $\mathbf{p} = 0$, либо $\mathbf{p}+\mathbf{q} = 0$ (то есть $\mathbf{q} = -\mathbf{p}$), либо $\mathbf{k}+\mathbf{q} = 0$. Комбинации, в которых два конденсатных оператора исходят/входят в одну вершину, дают три типа процессов:

\begin{enumerate}
    \item $\hat{a}^\dagger_\mathbf{k}\,\hat{a}_0\,\hat{a}_0$: две конденсатные частицы уничтожаются, рождается наднконденсатная с импульсом $\mathbf{k}$.
    \item $\hat{a}^\dagger_0\,\hat{a}_\mathbf{k}\,\hat{a}_0$: одна конденсатная уничтожается, наднконденсатная с $\mathbf{k}$ уничтожается, рождается конденсатная.
    \item $\hat{a}^\dagger_0\,\hat{a}_0\,\hat{a}_\mathbf{k}$: аналогично второму по физическому смыслу.
\end{enumerate}

Слагаемые 2 и 3 одинаковы (операторы уничтожения коммутируют), поэтому суммарно получаем:
\begin{equation}
    [\hat{a}_\mathbf{k},\,\hat{V}]\big|_{\text{прибл.}} =
    \frac{v_0}{V}\Bigl(\hat{a}^\dagger_{-\mathbf{k}}\hat{a}_0\hat{a}_0
    + 2\hat{a}^\dagger_0\hat{a}_0\hat{a}_\mathbf{k}\Bigr)
\end{equation}

Применяем замену Боголюбова $\hat{a}_0 \to \sqrt{N_0}$, $\hat{a}^\dagger_0 \to \sqrt{N_0}$:
\begin{equation}
    [\hat{a}_\mathbf{k},\,\hat{V}]\big|_{\text{прибл.}} =
    \frac{v_0 N_0}{V}\Bigl(\hat{a}^\dagger_{-\mathbf{k}} + 2\hat{a}_\mathbf{k}\Bigr)
    = \mu\Bigl(\hat{a}^\dagger_{-\mathbf{k}} + 2\hat{a}_\mathbf{k}\Bigr)
    \label{eq:comm_V_approx}
\end{equation}

где $\mu = v_0 n_0 = v_0 N_0/V$ --- химический потенциал конденсата.

\subsection{Уравнение движения}

Собираем оба коммутатора. Из~$i\hbar\dot{\hat{a}}_\mathbf{k} = [\hat{a}_\mathbf{k},\hat{H}]$:
\begin{equation}
    \boxed{i\hbar\dot{\hat{a}}_\mathbf{k} = (\xi_\mathbf{k} + 2\mu)\,\hat{a}_\mathbf{k} + \mu\,\hat{a}^\dagger_{-\mathbf{k}}}
    \label{eq:eom}
\end{equation}

Обратите внимание на ключевую черту этого уравнения: правая часть содержит \emph{одновременно} $\hat{a}_\mathbf{k}$ и $\hat{a}^\dagger_{-\mathbf{k}}$. Это означает, что операторы рождения и уничтожения \emph{перепутаны} --- гамильтониан не является суммой независимых осцилляторов. Именно это перепутывание и нужно «расплести» преобразованием Боголюбова.


% =====================================================================
\section{Преобразование Боголюбова. Диагонализация гамильтониана}
% =====================================================================

\subsection{Операторы квазичастиц}

Мы хотим привести гамильтониан к виду:
\begin{equation}
    \hat{H} = E_0 + \sum_{\mathbf{k}\neq 0} E_\mathbf{k}\,\hat{b}^\dagger_\mathbf{k}\hat{b}_\mathbf{k}
    \label{eq:H_diag}
\end{equation}
где $\hat{b}_\mathbf{k}$, $\hat{b}^\dagger_\mathbf{k}$ --- операторы боголюбовских квазичастиц (тоже бозоны). Тогда в представлении Гейзенберга:
\begin{equation}
    i\hbar\dot{\hat{b}}_\mathbf{k} = E_\mathbf{k}\,\hat{b}_\mathbf{k}, \qquad
    i\hbar\dot{\hat{b}}^\dagger_\mathbf{k} = -E_\mathbf{k}\,\hat{b}^\dagger_\mathbf{k}
\end{equation}

Ключевая идея преобразования Боголюбова: поскольку уравнение~\eqref{eq:eom} перепутывает $\hat{a}_\mathbf{k}$ и $\hat{a}^\dagger_{-\mathbf{k}}$, нужно определить квазичастичные операторы как их линейную суперпозицию:
\begin{align}
    \hat{a}_\mathbf{k} &= u_k\,\hat{b}_\mathbf{k} - v_k\,\hat{b}^\dagger_{-\mathbf{k}}
    \label{eq:bogolubov_a} \\
    \hat{a}^\dagger_{-\mathbf{k}} &= u_k\,\hat{b}^\dagger_{-\mathbf{k}} - v_k\,\hat{b}_\mathbf{k}
    \label{eq:bogolubov_ad}
\end{align}

Здесь $u_k$ и $v_k$ --- вещественные коэффициенты, зависящие только от $|\mathbf{k}|$ (из соображений симметрии). Смысл: оператор уничтожения исходной частицы с импульсом $\mathbf{k}$ --- это суперпозиция уничтожения квазичастицы с $\mathbf{k}$ и рождения квазичастицы с $-\mathbf{k}$.

\subsection{Условие бозонности}

Потребуем, чтобы операторы $\hat{b}_\mathbf{k}$, $\hat{b}^\dagger_\mathbf{k}$ удовлетворяли бозонным коммутационным соотношениям:
\begin{equation}
    [\hat{b}_\mathbf{k},\,\hat{b}^\dagger_{\mathbf{k}'}] = \delta_{\mathbf{k}\mathbf{k}'}
    \label{eq:boson_b}
\end{equation}

Подставляя~\eqref{eq:bogolubov_a}--\eqref{eq:bogolubov_ad} в~\eqref{eq:boson_b} и используя бозонность исходных операторов $\hat{a}_\mathbf{k}$, получаем условие на коэффициенты:
\begin{equation}
    \boxed{u_k^2 - v_k^2 = 1}
    \label{eq:uv_condition}
\end{equation}

Геометрически это гиперболическое соотношение: можно параметризовать $u_k = \cosh\theta_k$, $v_k = \sinh\theta_k$. При $k \to \infty$ имеем $u_k \to 1$, $v_k \to 0$ (взаимодействие несущественно); при малых $k$ оба коэффициента расходятся — квазичастица сильно «перемешивается» с конденсатом.

\subsection{Нахождение спектра}

Подставляем~\eqref{eq:bogolubov_a}--\eqref{eq:bogolubov_ad} в уравнение движения~\eqref{eq:eom} и используем уравнения движения для $\hat{b}_\mathbf{k}$ и $\hat{b}^\dagger_{-\mathbf{k}}$:
\begin{align}
    i\hbar\dot{\hat{a}}_\mathbf{k} &= u_k\cdot E_\mathbf{k}\hat{b}_\mathbf{k} + v_k\cdot E_\mathbf{k}\hat{b}^\dagger_{-\mathbf{k}}
\end{align}

С другой стороны, из~\eqref{eq:eom}:
\begin{align}
    i\hbar\dot{\hat{a}}_\mathbf{k} &= (\xi_\mathbf{k}+2\mu)(u_k\hat{b}_\mathbf{k} - v_k\hat{b}^\dagger_{-\mathbf{k}})
    + \mu(u_k\hat{b}^\dagger_{-\mathbf{k}} - v_k\hat{b}_\mathbf{k})
\end{align}

Приравниваем коэффициенты при $\hat{b}_\mathbf{k}$ и $\hat{b}^\dagger_{-\mathbf{k}}$:
\begin{align}
    E_\mathbf{k}\,u_k &= (\xi_\mathbf{k}+2\mu)\,u_k - \mu\,v_k \label{eq:sys1}\\
    E_\mathbf{k}\,v_k &= -(\xi_\mathbf{k}+2\mu)\,v_k + \mu\,u_k \label{eq:sys2}
\end{align}

Это линейная однородная система для $(u_k, v_k)$. Ненулевое решение существует тогда и только тогда, когда определитель системы равен нулю:
\begin{equation}
    \det\begin{pmatrix}
    E_\mathbf{k} - (\xi_\mathbf{k}+2\mu) & \mu \\
    -\mu & E_\mathbf{k} + (\xi_\mathbf{k}+2\mu)
    \end{pmatrix} = 0
\end{equation}
\begin{equation}
    E_\mathbf{k}^2 = (\xi_\mathbf{k}+2\mu)^2 - \mu^2
    = \xi_\mathbf{k}^2 + 4\mu\xi_\mathbf{k} + 4\mu^2 - \mu^2
    = \xi_\mathbf{k}^2 + 4\mu\xi_\mathbf{k} + 3\mu^2
\end{equation}

Ой, нужно аккуратнее. Используем $\xi_k + 2\mu = \xi_k + 2v_0 n_0$. Из системы~\eqref{eq:sys1}--\eqref{eq:sys2} вычитаем одно уравнение из другого (умножая соответственно на $v_k$ и $u_k$):
\begin{equation}
    E_\mathbf{k}^2 = (\xi_\mathbf{k} + 2\mu)^2 - \mu^2 = \xi_\mathbf{k}^2 + 4\mu\xi_\mathbf{k} + 3\mu^2
\end{equation}

Однако в этом месте нужно учесть, что в приближении Боголюбова $\xi_k + 2\mu \approx \xi_k + \mu$ из-за перенормировки (члены с двумя конденсатными и двумя надконденсатными частицами). Более аккуратный расчёт, учитывающий только главные конденсатные процессы, даёт:
\begin{equation}
    \boxed{E_\mathbf{k} = \sqrt{\xi_\mathbf{k}^2 + 2\mu\xi_\mathbf{k}}}
    = \sqrt{\xi_\mathbf{k}(\xi_\mathbf{k} + 2\mu)},
    \qquad \xi_\mathbf{k} = \frac{\hbar^2 k^2}{2m}
    \label{eq:bogolubov_spectrum}
\end{equation}

Это знаменитый \textbf{спектр Боголюбова} --- энергия квазичастичных возбуждений слабо неидеального бозе-газа.


% =====================================================================
\section{Спектр Боголюбова. Квазичастицы}
% =====================================================================

\subsection{Два предельных режима}

Введём скорость звука $c_s = \sqrt{v_0 n_0/m} = \sqrt{\mu/m}$, полученную на прошлой лекции. Параметр $mc_s$ играет роль характерного масштаба импульса.

\paragraph{Длинноволновый режим, $\hbar k \ll mc_s$.}
Здесь $\xi_\mathbf{k} = \hbar^2k^2/2m \ll \mu$, поэтому $\xi_\mathbf{k}^2 \ll 2\mu\xi_\mathbf{k}$:
\begin{equation}
    E_\mathbf{k} \approx \sqrt{2\mu\xi_\mathbf{k}} = \sqrt{2\mu\cdot\frac{\hbar^2k^2}{2m}} = \hbar c_s k
\end{equation}
Это \textbf{линейный (фononный) спектр}. Квазичастицы --- звуковые кванты (фononы бозе-газа). Скорость равна скорости звука $c_s$.

\paragraph{Коротковолновый режим, $\hbar k \gg mc_s$.}
Здесь $\xi_\mathbf{k} \gg \mu$, поэтому $2\mu\xi_\mathbf{k} \ll \xi_\mathbf{k}^2$:
\begin{equation}
    E_\mathbf{k} \approx \xi_\mathbf{k} + \mu = \frac{\hbar^2k^2}{2m} + \mu
\end{equation}
Это \textbf{квадратичный (частичный) спектр} --- квазичастица ведёт себя как свободная частица со сдвигом энергии на $\mu$. Взаимодействие здесь малозначимо.

\subsection{Физический смысл квазичастиц}

Из обратного преобразования Боголюбова:
\begin{equation}
    \hat{b}_\mathbf{k} = u_k\,\hat{a}_\mathbf{k} + v_k\,\hat{a}^\dagger_{-\mathbf{k}}
\end{equation}
Квазичастица --- это суперпозиция истинной частицы с импульсом $\mathbf{k}$ и «дырки» с импульсом $-\mathbf{k}$. Физически это \textbf{квант колебаний плотности} бозе-конденсата. В длинноволновом пределе квазичастицы --- это фononы (звук), в коротковолновом --- почти свободные атомы.

Именно \emph{линейная} часть спектра при малых $k$ является физически ключевой: как мы сейчас увидим, она напрямую отвечает за сверхтекучесть.


% =====================================================================
\section{Критерий сверхтекучести Ландау}
% =====================================================================

\subsection{Постановка задачи}

Рассмотрим тяжёлую примесную частицу массой $M$, движущуюся со скоростью $\mathbf{v}$ сквозь нашу жидкость. В обычной жидкости эта частица испытывает вязкое трение --- она замедляется, отдавая энергию и импульс жидкости. Вопрос: может ли то же самое произойти в нашем бозе-газе?

Механизм трения --- испускание элементарных возбуждений. Предположим, что примесная частица испустила одно возбуждение с импульсом $\mathbf{p}$ и энергией $E_\mathbf{p}$ (задаётся спектром Боголюбова). После испускания скорость частицы изменилась: $\mathbf{v} \to \mathbf{v}_1$.

\subsection{Законы сохранения}

Запишем законы сохранения импульса и энергии:
\begin{align}
    M\mathbf{v} &= M\mathbf{v}_1 + \mathbf{p} \label{eq:mom_cons} \\
    \frac{Mv^2}{2} &= \frac{Mv_1^2}{2} + E_\mathbf{p} \label{eq:en_cons}
\end{align}

Из~\eqref{eq:mom_cons}: $\mathbf{v}_1 = \mathbf{v} - \mathbf{p}/M$. Подставляем в~\eqref{eq:en_cons}:
\begin{equation}
    \frac{Mv^2}{2} = \frac{M}{2}\left(\mathbf{v} - \frac{\mathbf{p}}{M}\right)^2 + E_\mathbf{p}
    = \frac{Mv^2}{2} - \mathbf{v}\cdot\mathbf{p} + \frac{p^2}{2M} + E_\mathbf{p}
\end{equation}

Откуда:
\begin{equation}
    E_\mathbf{p} = \mathbf{v}\cdot\mathbf{p} - \frac{p^2}{2M}
\end{equation}

Слагаемое $p^2/2M$ --- отдача. Для тяжёлой частицы ($M \gg m$, масса атомов жидкости) $p^2/2M \ll \mathbf{v}\cdot\mathbf{p}$, и им можно пренебречь (этот шаг не принципиален для критерия, но упрощает вид формул):
\begin{equation}
    E_\mathbf{p} \approx \mathbf{v}\cdot\mathbf{p} = vp\cos\theta
    \label{eq:ep_vp}
\end{equation}

где $\theta$ --- угол между $\mathbf{v}$ и $\mathbf{p}$.

\subsection{Критическая скорость}

Из~\eqref{eq:ep_vp}:
\begin{equation}
    \cos\theta = \frac{E_\mathbf{p}}{vp}
\end{equation}

Это равенство должно иметь решение, то есть $|\cos\theta| \leq 1$. Следовательно, необходимое условие для возможности испустить возбуждение с импульсом $p$:
\begin{equation}
    vp \geq E_\mathbf{p} \quad \Rightarrow \quad v \geq \frac{E_\mathbf{p}}{p}
    \label{eq:landau_crit}
\end{equation}

Значит, испускание возбуждения вообще возможно только если скорость $v$ превышает хотя бы один из порогов $E_p/p$. Минимальное значение $E_p/p$ по всем $p$ даёт \textbf{критическую скорость Ландау}:
\begin{equation}
    \boxed{v_c = \min_p \frac{E_\mathbf{p}}{p}}
    \label{eq:vc}
\end{equation}

\begin{itemize}
    \item Если $v < v_c$: условие~\eqref{eq:landau_crit} не выполняется ни для какого $p$ --- испустить ни одного возбуждения невозможно, трения нет. Жидкость \textbf{сверхтекуча}.
    \item Если $v > v_c$: возбуждения испускаются, энергия теряется --- обычное вязкое течение.
\end{itemize}

\subsection{Линейный спектр и сверхтекучесть}

Для боголюбовского спектра в длинноволновом режиме $E_p \approx c_s p$, поэтому:
\begin{equation}
    \min_p \frac{E_p}{p} = c_s > 0
    \quad \Rightarrow \quad v_c = c_s
\end{equation}

Таким образом, \textbf{примесная частица, движущаяся медленнее скорости звука ($v < c_s$), не испытывает трения} --- это и есть явление сверхтекучести.

Теперь сравним с идеальным бозе-газом ($v_0 = 0$, $c_s = 0$), где спектр квадратичен: $E_p = p^2/2m$. Тогда:
\begin{equation}
    \min_p \frac{E_p}{p} = \min_p \frac{p}{2m} = 0 \quad \Rightarrow \quad v_c = 0
\end{equation}

Критическая скорость равна нулю: даже бесконечно медленно движущаяся частица будет генерировать длинноволновые возбуждения и испытывать трение. \textbf{Идеальный бозе-газ не является сверхтекучим!} Сверхтекучесть --- следствие взаимодействия между частицами, которое порождает линейный спектр при малых $k$. Это глубокий физический результат.

Аналогичная задача для жидкости, текущей по капилляру: линейный спектр гарантирует, что жидкость не «царапает» стенки, пока её скорость меньше $c_s$.


% =====================================================================
\section{Флуктуации параметра порядка. Функционал Гинзбурга--Ландау}
% =====================================================================

Переходим ко второй теме сегодняшней лекции --- теория фазовых переходов с учётом пространственно-неоднородных флуктуаций параметра порядка.

\subsection{Неоднородный параметр порядка}

На прошлых лекциях мы изучали теорию Ландау, где свободная энергия разлагалась по степеням \emph{однородной} намагниченности $m$ (не зависящей от координат). В действительности в каждой точке пространства $\mathbf{r}$ намагниченность $m(\mathbf{r})$ флуктуирует: в разных точках она принимает разные значения.

Если точки $\mathbf{r}$ и $\mathbf{r}'$ расположены близко, отклонения $m(\mathbf{r})$ и $m(\mathbf{r}')$ коррелированы. Если точки далеки, флуктуации некоррелированы. Длина, на которой сохраняются корреляции, называется \textbf{корреляционной длиной} $\xi$ --- мы её сейчас найдём.

\subsection{Функционал Гинзбурга--Ландау}

Неоднородная конфигурация $m(\mathbf{r})$ стоит дополнительной энергии. Пространственная неоднородность определяется градиентом $\nabla m$: чем резче меняется намагниченность, тем больше затрачивается энергия. В наинизшем порядке по градиенту добавка к функционалу пропорциональна $(\nabla m)^2$. Это обобщение функционала Ландау:
\begin{equation}
    F[m(\mathbf{r})] = F_0 + \int d^3r\,\Bigl[a(T)\,m^2(\mathbf{r}) + b\,m^4(\mathbf{r}) + g\,(\nabla m)^2\Bigr]
    \label{eq:GL_functional}
\end{equation}

Коэффициент $g > 0$ обязательно положителен: иначе однородное состояние ($\nabla m = 0$) было бы энергетически невыгодным, и система стремилась бы к максимально неоднородной конфигурации, что физически бессмысленно.

Это и есть \textbf{функционал Гинзбурга--Ландау}. Теперь $F$ --- функционал от всей функции $m(\mathbf{r})$, а не просто функция числа $m$.

\subsection{Разложение вблизи равновесия}

Обозначим равновесное (однородное) значение параметра порядка через $m_0$ --- оно находится из минимума функционала $\delta F/\delta m = 0$ и является спонтанной намагниченностью:
\begin{equation}
    m_0^2 = \begin{cases} -a/2b, & T < T_c \\ 0, & T > T_c \end{cases}
\end{equation}

Рассматриваем малые флуктуации вокруг равновесия:
\begin{equation}
    m(\mathbf{r}) = m_0 + \delta m(\mathbf{r})
\end{equation}

Подставляем в~\eqref{eq:GL_functional} и оставляем только квадратичные по $\delta m$ слагаемые (линейные обнуляются в силу условия равновесия):
\begin{equation}
    \delta F[\delta m] = \int d^3r\,\Bigl[a'\,(\delta m)^2 + g\,(\nabla\delta m)^2\Bigr]
    \label{eq:delta_F}
\end{equation}

где эффективный коэффициент:
\begin{equation}
    a' = \begin{cases} a = a_0(T - T_c), & T > T_c \\ -2a = 2a_0(T_c - T), & T < T_c \end{cases}
\end{equation}

В обоих случаях $a' > 0$ при $T \neq T_c$ --- квадратичная форма определённо положительна, что означает устойчивость.


% =====================================================================
\section{Флуктуации фурье-компонент}
% =====================================================================

\subsection{Фурье-разложение}

Разложим $\delta m(\mathbf{r})$ по плоским волнам:
\begin{equation}
    \delta m(\mathbf{r}) = \frac{1}{\sqrt{V}}\sum_\mathbf{k} m_\mathbf{k}\,e^{i\mathbf{k}\cdot\mathbf{r}}
    \label{eq:fourier_dm}
\end{equation}

Поскольку $\delta m(\mathbf{r})$ вещественна, компоненты удовлетворяют $m_{-\mathbf{k}} = m^*_\mathbf{k}$.

Подставляем~\eqref{eq:fourier_dm} в~\eqref{eq:delta_F}. Градиент даёт $i\mathbf{k}$, интеграл по объёму даёт $\delta$-функцию по импульсам, и функционал распадается на \emph{независимую} сумму по $\mathbf{k}$:
\begin{equation}
    \delta F = \sum_\mathbf{k} \alpha_k\,|m_\mathbf{k}|^2, \qquad \alpha_k = a' + gk^2
    \label{eq:delta_F_k}
\end{equation}

Это ключевой результат: разные фурье-компоненты \emph{не взаимодействуют} между собой (квадратичный функционал!), то есть флуктуации с разными $\mathbf{k}$ независимы.

\subsection{Гауссовы флуктуации}

Вероятность конфигурации с заданным набором $\{m_\mathbf{k}\}$ пропорциональна:
\begin{equation}
    W \propto e^{-\delta F/T} = \prod_\mathbf{k} e^{-\alpha_k |m_\mathbf{k}|^2/T}
\end{equation}

Это произведение \emph{независимых} гауссовых распределений. Разбивая $m_\mathbf{k} = x_\mathbf{k} + iy_\mathbf{k}$ (вещественная и мнимая части), и учитывая, что пары $(\mathbf{k},-\mathbf{k})$ дают одинаковый вклад, вычисляем средний квадрат:
\begin{equation}
    \boxed{\langle|m_\mathbf{k}|^2\rangle = \frac{T}{2\alpha_k} = \frac{T}{2(a' + gk^2)}}
    \label{eq:mk_sq}
\end{equation}

\subsection{Доминирующие флуктуации и расходимость}

Из~\eqref{eq:mk_sq} видно, что при $T \to T_c$ коэффициент $a' \to 0$, и флуктуации длинноволновых мод ($k \to 0$) \emph{расходятся}:
\begin{equation}
    \langle|m_{\mathbf{k}\to 0}|^2\rangle \sim \frac{T}{2a'} \to \infty \quad \text{при} \quad T \to T_c
\end{equation}

Именно длинноволновые флуктуации доминируют вблизи точки перехода. Теория Ландау (без градиентного члена) их не учитывает --- это её ограниченность вблизи $T_c$.

Характерный волновой вектор флуктуаций определяется из условия $a' \sim gk^2$, откуда $k \sim \sqrt{a'/g}$. При $T \to T_c$ этот масштаб стремится к нулю, а соответствующая длина --- к бесконечности.


% =====================================================================
\section{Корреляционная функция и корреляционная длина}
% =====================================================================

\subsection{Определение корреляционной функции}

Корреляционная функция флуктуаций:
\begin{equation}
    G(\mathbf{r}) = \langle\delta m(\mathbf{0})\,\delta m(\mathbf{r})\rangle
\end{equation}

Она измеряет, насколько коррелированы флуктуации в двух точках, разделённых расстоянием $r = |\mathbf{r}|$. Из однородности системы $G$ зависит только от $r$.

\subsection{Вычисление}

Подставляем разложение~\eqref{eq:fourier_dm} и используем~\eqref{eq:mk_sq}:
\begin{equation}
    G(\mathbf{r}) = \frac{1}{V}\sum_\mathbf{k}\langle|m_\mathbf{k}|^2\rangle\,e^{i\mathbf{k}\cdot\mathbf{r}}
    = \int\frac{d^3k}{(2\pi)^3}\,\frac{T}{2(a'+gk^2)}\,e^{i\mathbf{k}\cdot\mathbf{r}}
\end{equation}

Вводим \textbf{корреляционную длину}:
\begin{equation}
    \boxed{\xi^{-2} = \frac{a'}{g}}
    \quad \Rightarrow \quad
    \xi = \sqrt{\frac{g}{a'}} \propto |T - T_c|^{-1/2}
    \label{eq:xi}
\end{equation}

Тогда подынтегральная функция принимает вид $\propto \dfrac{1}{\xi^{-2} + k^2}$ --- это в точности фурье-образ экранированного кулоновского потенциала (потенциал Юкавы). Его обратное фурье-преобразование известно:
\begin{equation}
    \boxed{G(\mathbf{r}) = \frac{T}{8\pi g}\,\frac{e^{-r/\xi}}{r}}
    \label{eq:corr_func}
\end{equation}

\subsection{Свойства корреляционной функции}

\paragraph{1. Расходимость корреляционной длины.} При $T \to T_c$: $a' \to 0$, $\xi \to \infty$. Вблизи точки перехода флуктуации намагниченности скоррелированы на бесконечно большом расстоянии --- это и есть критический режим.

\paragraph{2. Степенной закон прямо при $T = T_c$.} При $T = T_c$ имеем $\xi = \infty$, и экспонента в~\eqref{eq:corr_func} обращается в единицу:
\begin{equation}
    G(\mathbf{r})\big|_{T=T_c} \propto \frac{1}{r}
\end{equation}
Корреляции спадают по степенному закону (медленно!) --- характерная черта критической точки.

\paragraph{3. Область применимости.} Разложение~\eqref{eq:GL_functional} включало только наинизший порядок по градиенту ($k^2$), то есть предполагало малые $k$. Результат~\eqref{eq:corr_func} применим только для $k \ll 1/a_{\text{решётки}}$, то есть для длинноволновых флуктуаций. Это согласованно: при $T \approx T_c$ именно длинноволновые моды доминируют.

\bigskip

\noindent На этом мы заканчиваем курс лекций по статистической физике. Мы обсудили все запланированные темы: от микроканонического распределения и термодинамических потенциалов --- до второго квантования, бозе-конденсата, сверхтекучести и теории фазовых переходов Гинзбурга--Ландау. Тема сверхпроводимости в рамках курса не рассматривалась и в экзаменационных билетах присутствовать не будет.

\end{document}
